\documentclass{article}
\newcommand{\mydate}{March 5, 2026}
\newcommand{\mytitle}{GP2 HW1}
\title{\textbf{\mytitle}}
\author{Jiete XUE\thanks{ID: 20253121013}}
\date{\mydate}
\usepackage{fancyhdr}
\pagestyle{fancy}
\fancyhf{}
\fancyhead[C]{\mytitle }
\fancyhead[R]{Jiete Xue}
\fancyhead[L]{\mydate}
\fancyfoot[C]{\thepage}
\usepackage{amsthm}
\usepackage{amsmath}
\usepackage{amssymb}
\usepackage{bm}
\usepackage{enumitem}
\usepackage{physics}
\usepackage{tikz}
\usetikzlibrary{arrows.meta, decorations.markings}
\usepackage{float}

%% 右矢
%\ket{\psi}          % 输出：|ψ⟩
%\ket{\psi(t)}       % 输出：|ψ(t)⟩
%
%% 左矢
%\bra{\phi}          % 输出：⟨φ|
%
%% 期望值
%\expval{\hat{A}}    % 输出：⟨Â⟩
%\expval{\hat{A}}{\psi}  % 输出：⟨ψ|Â|ψ⟩
%
%% 对易子
%\comm{\hat{A}}{\hat{B}}  % 输出：[Â, B̂]
\newtheoremstyle{1}{}{}{}{}{\bfseries}{}{\newline}{}
\theoremstyle{1}
\newtheorem{problem}{Problem}
\newtheorem{solution}{Solution}
\usepackage{chngcntr}
\counterwithin{equation}{problem}

\setlist[enumerate]{label=(\arabic*), leftmargin=*, align=left}

\newcommand{\pa}{\partial}
\newcommand{\ii}{\mathrm{i}}
\newcommand{\ee}{\mathrm{e}}

\begin{document}
\maketitle
\begin{center}
    {\Large\textbf{Why these boring maths…}}
\end{center}
\begin{problem}[Levi-Civita symbol, Kronecker symbol and repeated
symbol represents summation]


\hspace{12pt} $\varepsilon_{ijk}$, $\varepsilon_{jki}$, $\varepsilon_{kij}$ have a relation of even permutations, by definition, they should have the same value.

One has
\begin{equation}
    \varepsilon^{ijk} \varepsilon_{lmn} = \delta^{\left[i\right.}_{l}\delta^{j}_{m}\delta^{\left.k\right]}_{n}.
\end{equation}

When $i=l$, that is
\begin{equation}
    \delta^{\left[i\right.}_{i}\delta^{j}_{m}\delta^{\left.k\right]}_{n} = (3-1)! \delta^{\left[j\right.}_{m}\delta^{\left.k\right]}_{n} = \delta^j_m \delta^k_n - \delta^k_m \delta^j_n.
\end{equation}
By substituting $(i,j,k)$ with $(k,i,j)$, we obtain another form.
\end{problem}
\begin{problem}[Vector algebra]
    \quad
    \begin{enumerate}[label=\arabic*.]
        \item 

    \begin{enumerate}
\item \(\vec{A} \cdot (\vec{B} \times \vec{C}) = \vec{B} \cdot (\vec{C} \times \vec{A}) = \vec{C} \cdot (\vec{A} \times \vec{B})\)

Left-hand side:
\begin{equation}
\vec{A} \cdot (\vec{B} \times \vec{C}) = A_i (\vec{B} \times \vec{C})_i
\end{equation}
\begin{equation}
(\vec{B} \times \vec{C})_i = \varepsilon_{ijk} B_j C_k
\end{equation}
Thus,
\begin{equation}
\vec{A} \cdot (\vec{B} \times \vec{C}) = A_i \varepsilon_{ijk} B_j C_k = \varepsilon_{ijk} A_i B_j C_k
\end{equation}

Similarly,
\begin{equation}
\vec{B} \cdot (\vec{C} \times \vec{A}) = B_i (\vec{C} \times \vec{A})_i = B_i \varepsilon_{ijk} C_j A_k
\end{equation}
Rename indices: \(i \to j, j \to k, k \to i\):
\begin{equation}
= \varepsilon_{jki} B_j C_k A_i
\end{equation}
Since \(\varepsilon_{jki} = \varepsilon_{ijk}\) (cyclic permutation does not change the sign in 3D):
\begin{equation}
= \varepsilon_{ijk} A_i B_j C_k
\end{equation}
The third expression follows the same way.

\begin{equation}
\boxed{\vec{A} \cdot (\vec{B} \times \vec{C}) = \vec{B} \cdot (\vec{C} \times \vec{A}) = \vec{C} \cdot (\vec{A} \times \vec{B})}
\end{equation}

\item \(\vec{A} \times (\vec{B} \times \vec{C}) = \vec{B}(\vec{A} \cdot \vec{C}) - \vec{C}(\vec{A} \cdot \vec{B})\)

Let \(\vec{D} = \vec{B} \times \vec{C}\). Then
\begin{equation}
D_i = \varepsilon_{ijk} B_j C_k
\end{equation}
\begin{equation}
[\vec{A} \times \vec{D}]_m = \varepsilon_{mni} A_n D_i
\end{equation}
Substitute \(D_i\):
\begin{equation}
[\vec{A} \times (\vec{B} \times \vec{C})]_m = \varepsilon_{mni} A_n \varepsilon_{ijk} B_j C_k
\end{equation}
Use \(\varepsilon_{mni} = \varepsilon_{imn}\) and the identity \(\varepsilon_{imn} \varepsilon_{ijk} = \delta_{mj}\delta_{nk} - \delta_{mk}\delta_{nj}\):
\begin{equation}
= (\delta_{mj}\delta_{nk} - \delta_{mk}\delta_{nj}) A_n B_j C_k
\end{equation}
Expand:
\begin{equation}
= \delta_{mj}\delta_{nk} A_n B_j C_k - \delta_{mk}\delta_{nj} A_n B_j C_k
\end{equation}
First term: \(\delta_{mj}\delta_{nk} A_n B_j C_k\):  
\(\delta_{mj} B_j = B_m\), \(\delta_{nk} C_k = C_n\), and \(A_n C_n = \vec{A} \cdot \vec{C}\).  
So first term = \(B_m (\vec{A} \cdot \vec{C})\).

Second term: \(\delta_{mk}\delta_{nj} A_n B_j C_k\):  
\(\delta_{mk} C_k = C_m\), \(\delta_{nj} B_j = B_n\), and \(A_n B_n = \vec{A} \cdot \vec{B}\).  
So second term = \(C_m (\vec{A} \cdot \vec{B})\).

Thus,
\begin{equation}
[\vec{A} \times (\vec{B} \times \vec{C})]_m = B_m (\vec{A} \cdot \vec{C}) - C_m (\vec{A} \cdot \vec{B})
\end{equation}

\begin{equation}
\boxed{\vec{A} \times (\vec{B} \times \vec{C}) = \vec{B}(\vec{A} \cdot \vec{C}) - \vec{C}(\vec{A} \cdot \vec{B})}
\end{equation}

\item \((\vec{A} \times \vec{B}) \times \vec{C} = -\vec{A}(\vec{B} \cdot \vec{C}) + \vec{B}(\vec{A} \cdot \vec{C})\)

Let \(\vec{D} = \vec{A} \times \vec{B}\). Then
\begin{equation}
D_k = \varepsilon_{kij} A_i B_j
\end{equation}
\begin{equation}
[\vec{D} \times \vec{C}]_m = \varepsilon_{mkl} D_k C_l
\end{equation}
Substitute \(D_k\):
\begin{equation}
[(\vec{A} \times \vec{B}) \times \vec{C}]_m = \varepsilon_{mkl} \varepsilon_{kij} A_i B_j C_l
\end{equation}
Use \(\varepsilon_{mkl} = -\varepsilon_{kml}\), so:
\begin{equation}
= -\varepsilon_{kml} \varepsilon_{kij} A_i B_j C_l
\end{equation}
Apply \(\varepsilon_{kml} \varepsilon_{kij} = \delta_{mi}\delta_{lj} - \delta_{mj}\delta_{li}\):
\begin{equation}
= - (\delta_{mi}\delta_{lj} - \delta_{mj}\delta_{li}) A_i B_j C_l
\end{equation}
\begin{equation}
= -\delta_{mi}\delta_{lj} A_i B_j C_l + \delta_{mj}\delta_{li} A_i B_j C_l
\end{equation}
First term: \(- A_m B_l C_l = -A_m (\vec{B} \cdot \vec{C})\)  
Second term: \(+ B_m A_l C_l = +B_m (\vec{A} \cdot \vec{C})\)

Thus:
\begin{equation}
[(\vec{A} \times \vec{B}) \times \vec{C}]_m = B_m (\vec{A} \cdot \vec{C}) - A_m (\vec{B} \cdot \vec{C})
\end{equation}
Therefore:
\begin{equation}
\boxed{(\vec{A} \times \vec{B}) \times \vec{C} = \vec{B}(\vec{A} \cdot \vec{C}) - \vec{A}(\vec{B} \cdot \vec{C})}
\end{equation}

\item \((\vec{A} \times \vec{B}) \cdot (\vec{C} \times \vec{D}) = (\vec{A} \cdot \vec{C})(\vec{B} \cdot \vec{D}) - (\vec{A} \cdot \vec{D})(\vec{B} \cdot \vec{C})\)

Let \(\vec{P} = \vec{A} \times \vec{B}\), \(\vec{Q} = \vec{C} \times \vec{D}\).

Then
\begin{equation}
P_i = \varepsilon_{ijk} A_j B_k,\quad Q_i = \varepsilon_{ilm} C_l D_m
\end{equation}
\begin{equation}
\vec{P} \cdot \vec{Q} = P_i Q_i = \varepsilon_{ijk} A_j B_k \cdot \varepsilon_{ilm} C_l D_m
\end{equation}
\begin{equation}
= \varepsilon_{ijk} \varepsilon_{ilm} A_j B_k C_l D_m
\end{equation}
Using \(\varepsilon_{ijk} \varepsilon_{ilm} = \delta_{jl}\delta_{km} - \delta_{jm}\delta_{kl}\):
\begin{equation}
= (\delta_{jl}\delta_{km} - \delta_{jm}\delta_{kl}) A_j B_k C_l D_m
\end{equation}
First term: \(\delta_{jl}\delta_{km} A_j B_k C_l D_m = A_j C_j \cdot B_k D_k = (\vec{A} \cdot \vec{C})(\vec{B} \cdot \vec{D})\).

Second term: \(-\delta_{jm}\delta_{kl} A_j B_k C_l D_m = - A_j D_j \cdot B_k C_k = -(\vec{A} \cdot \vec{D})(\vec{B} \cdot \vec{C})\).

Thus:
\begin{equation}
\boxed{(\vec{A} \times \vec{B}) \cdot (\vec{C} \times \vec{D}) = (\vec{A} \cdot \vec{C})(\vec{B} \cdot \vec{D}) - (\vec{A} \cdot \vec{D})(\vec{B} \cdot \vec{C})}
\end{equation}
    \end{enumerate}
\item \begin{enumerate}
    \item \textbf{Similarly:} Both transform like vectors under proper rotations (determinant = +1).
    \item \textbf{Difference:} Under spatial inversion, Polar vector changes sign, but pseudovector does not change sign.
    \item \textbf{Examples of pseudovectors:} Angular momentum, Torque.
\end{enumerate}
\end{enumerate}
\end{problem}
\begin{problem}[The vector calculus]
    \quad
    \begin{enumerate}
        \item \begin{equation}
            \pa_i(fg) = \pa_i f g + f\pa_i g.
        \end{equation}
        \begin{equation}
            \pa_i\left(\frac{f}{g}\right) = \pa_i f \frac{1}{g} + f \pa_i \left(\frac{1}{g}\right) = \frac{\pa_i f}{g} - f\left( \frac{\pa_i g}{g^2}\right).
        \end{equation}
        \item \begin{equation}
            \pa_i \left(f A_i\right) =  f\left(\pa_i A_i\right)  + \pa_i f A_i = f \left(\nabla \cdot \vec{A}\right) + \nabla f \cdot\vec{A}.
        \end{equation}
        \begin{equation}
            \pa_i \left(f A_j\right) \varepsilon_{ijk} = \pa_i f A_j \varepsilon_{ijk} + f \pa_i A_j \varepsilon_{ijk}= \left[f\left(\nabla \times \vec{A}\right)  - \vec{A} \times\left(\nabla f\right)\right]_{k}.
        \end{equation}
        \begin{equation}
            \pa_k \left(A_i B_j \varepsilon_{ijk}\right) = \left(\pa_k A_i B_j +A_i \pa_k B_j\right) \varepsilon_{ijk} = \vec{B} \cdot\left(\nabla \times \vec{A}\right) - \vec{A} \cdot\left( \nabla \times \vec{B}\right).
        \end{equation}
        \item Similarly.
    \end{enumerate}
\end{problem}
\begin{problem}[Fundamental theorems of vector calculus]
    \begin{enumerate}
        \item Under the spherical system,
        \begin{equation}
            \nabla \cdot\left(\frac{q}{r^2}\hat{\mathbf{r}}\right) = \frac{\pa}{\pa r}\left(q\right) = 0.\ (r\neq 0)
        \end{equation}
        It doesn't violate Gauss's theorem.
        \item Since $\vec{E}$ is isotropic, $\nabla \times \vec{E}=0$. By Stokes' theorem, we have
        \begin{equation}
            \oint_C \vec{E}\cdot\dd{\vec{l}} = 0.
        \end{equation}
        Let
        \begin{equation}
            \phi = \frac{q}{r},
        \end{equation}
        then, $\vec{E} = -\nabla \phi$.
    \end{enumerate}
\end{problem}

\begin{problem}[Multi-valued functions]
    \quad
    \begin{enumerate}

        \item One has,
        \begin{equation}
            \dd{\mathbf{r}}\cdot\nabla \theta = \dd{\theta}
        \end{equation}
        So,
        \begin{equation}
            \oint \dd{\mathbf{r}}\cdot\nabla \theta = \oint \dd{\theta } =2\pi n,
        \end{equation}
        where $n$ is the number of circuit goes around the origin.
        \item By Stokes theorem, 
        \begin{equation}
            \oint \dd{\mathbf{r}} \cdot \nabla\theta = \iint \dd{\mathbf{S}}\cdot \nabla\times\nabla\theta.
        \end{equation}
        
        While left-hand side is a constant, right-hand side is range dependent, which lead to
        \begin{equation}
            \nabla\times\nabla\theta = 2\pi n\delta^2(0).
        \end{equation}
    \end{enumerate}
\end{problem}
\end{document}